\documentclass[11pt,a4paper]{amsart}

\usepackage{amsmath,amssymb,amsthm}
\usepackage{hyperref}
\usepackage{booktabs}
\usepackage{microtype}

% ── Theorem environments ─────────────────────────────────────────────────────
\newtheorem{theorem}{Theorem}
\newtheorem{corollary}[theorem]{Corollary}
\newtheorem{lemma}[theorem]{Lemma}
\newtheorem{remark}[theorem]{Remark}
\newtheorem{definition}[theorem]{Definition}

% ── Macros ───────────────────────────────────────────────────────────────────
\newcommand{\eml}{\mathrm{eml}}
\newcommand{\R}{\mathbb{R}}
\newcommand{\C}{\mathbb{C}}
\newcommand{\N}{\mathbb{N}}
\newcommand{\Z}{\mathbb{Z}}

% ── Metadata ─────────────────────────────────────────────────────────────────
\title{The Infinite Zeros Barrier:\\
  \large $\sin(x)$ Has No Finite Real EML Tree}

\author{Art Almaguer}
\address{Independent researcher, Seattle, WA}
\email{contact@monogate.dev}
\date{\today}

\keywords{EML operator, elementary functions, exp-minus-log, tree representation,
  infinite zeros, symbolic computation}

\subjclass[2020]{26A09, 68W30, 30D15}

\begin{document}

\begin{abstract}
The EML operator $\eml(x,y) = \exp(x) - \ln(y)$ generates all elementary
functions by finite composition from the constant~$1$
\cite{odzryzwolek2026eml}.
We prove that $\sin(x)$ is not representable as any finite real EML tree
under the strict grammar (principal-branch $\ln$, $\ln(0)$ undefined).
The proof uses two classical facts: finite compositions of $\exp$ and $\ln$
are real-analytic, and real-analytic functions that are not identically zero
have only isolated zeros.  Since $\sin(x)$ has infinitely many zeros on any
unbounded interval, no finite real EML tree can equal it.
We name this the \emph{Infinite Zeros Barrier}.
Exhaustive computational search over all real EML trees with up to~12
internal nodes ($\sim\!1.7 \times 10^9$ trees) confirms no tree approximates
$\sin(x)$ with mean squared error below $10^{-6}$.
\end{abstract}

\maketitle

% ─────────────────────────────────────────────────────────────────────────────
\section{Introduction}
% ─────────────────────────────────────────────────────────────────────────────

In~\cite{odzryzwolek2026eml}, Odrzywołek introduced the binary operator
\[
  \eml(x,y) \;=\; \exp(x) - \ln(y)
\]
and proved that every elementary function over the reals is a finite composition
of~$\eml$ applied to the single constant~$1$.  The \emph{EML depth} of a function
$f$ is the number of $\eml$ nodes in the smallest such tree.

A natural question arises immediately: which elementary functions have finite
real EML depth, and which do not?  Odrzywołek's result guarantees representability
over~$\C$ (complex EML), but the real case is more restrictive.

\begin{definition}[Real EML tree]
  A \emph{real EML tree} of depth~$k$ is a full binary tree with $k$ internal
  nodes, each labeled~$\eml$, and leaves labeled with elements of~$\R$.
  Under the \emph{strict grammar}, the constant leaf set is $\{1\}$, and
  $\ln(y)$ is defined only for $y > 0$ (principal branch; $\ln(0)$ is undefined).
  The \emph{evaluation} of such a tree at~$x$ substitutes $x$ for variable leaves
  and $1$ for constant leaves.
\end{definition}

\begin{definition}[Finite real EML representability]
  A function $f : \R \to \R$ is \emph{finitely real-EML representable} if
  there exists a real EML tree~$T$ such that $T(x) = f(x)$ for all~$x$ in
  the domain of~$f$.
\end{definition}

Our main result:

\begin{theorem}[Infinite Zeros Barrier]
  \label{thm:main}
  $\sin(x)$ is not finitely real-EML representable.
\end{theorem}

The same argument applies to $\cos(x)$, $\tan(x)$, and every real-valued
function with infinitely many zeros on an unbounded domain.

% ─────────────────────────────────────────────────────────────────────────────
\section{The Proof}
% ─────────────────────────────────────────────────────────────────────────────

The proof rests on two pillars: real-analyticity of EML trees, and the
identity theorem for real-analytic functions.

\begin{lemma}[EML trees are real-analytic]
  \label{lem:analytic}
  Let $T$ be a real EML tree that is defined on an open interval $I \subset \R$.
  Then $T|_I$ is real-analytic.
\end{lemma}

\begin{proof}
  By induction on the number of internal nodes~$k$.

  \emph{Base case ($k = 0$):} Every leaf is either the constant~$1$ or the
  variable~$x$, both of which are real-analytic.

  \emph{Inductive step:} Suppose both subtrees $L$ and $R$ define real-analytic
  functions $\ell$ and $r$ on an open interval~$I$.  Then
  \[
    T(x) \;=\; \eml(\ell(x),\, r(x)) \;=\; e^{\ell(x)} - \ln(r(x)).
  \]
  Since $e^z$ is entire and composition of real-analytic functions is real-analytic,
  $e^{\ell(x)}$ is real-analytic on~$I$.  Since $\ln$ is real-analytic on
  $(0,\infty)$ and $r(x) > 0$ on~$I$ (strict grammar), $\ln(r(x))$ is
  real-analytic on~$I$.  Their difference is real-analytic on~$I$. \qed
\end{proof}

\begin{lemma}[Isolated zeros]
  \label{lem:isolated}
  Let $f$ be real-analytic on a connected open set $U \subset \R$ and
  not identically zero.  Then the zeros of $f$ are isolated:
  for every $x_0 \in U$ with $f(x_0) = 0$ there exists $\delta > 0$
  such that $f(x) \neq 0$ for $0 < |x - x_0| < \delta$.
\end{lemma}

\begin{proof}
  Standard consequence of the identity theorem for real-analytic functions;
  see e.g.~\cite{krantz2002primer}~Theorem~1.1.5.
\end{proof}

\begin{proof}[Proof of Theorem~\ref{thm:main}]
  Suppose for contradiction that $T$ is a real EML tree with $T(x) = \sin(x)$
  for all $x$ in a connected open domain $U$.

  By Lemma~\ref{lem:analytic}, $T|_U$ is real-analytic.  By Lemma~\ref{lem:isolated},
  its zeros are isolated.

  However, $\sin(x)$ vanishes at $x = n\pi$ for every $n \in \Z$.
  Every bounded subinterval of $U$ contains infinitely many zeros of the form
  $n\pi$, which are not isolated (they accumulate at $\pm\infty$).
  More precisely, the sequence $\{n\pi\}_{n \in \Z} \cap U$ is infinite with
  no isolated point in any bounded subinterval of $U$ large enough to contain
  two consecutive values.

  This contradicts the isolated-zeros property.  Therefore no such tree~$T$
  exists. \qed
\end{proof}

% ─────────────────────────────────────────────────────────────────────────────
\section{Corollaries}
% ─────────────────────────────────────────────────────────────────────────────

\begin{corollary}
  $\cos(x)$ is not finitely real-EML representable.
\end{corollary}
\begin{proof}
  $\cos(x)$ has zeros at $x = (2n+1)\pi/2$ for $n \in \Z$, an infinite sequence.
  The same argument applies. \qed
\end{proof}

\begin{corollary}
  $\tan(x)$ is not finitely real-EML representable on any interval containing
  two consecutive zeros.
\end{corollary}
\begin{proof}
  $\tan(x) = 0$ at $x = n\pi$ for all $n \in \Z$. Same argument. \qed
\end{proof}

\begin{corollary}[General Infinite Zeros Barrier]
  \label{cor:general}
  Let $f : \R \to \R$ be a function that has infinitely many zeros.
  Then $f$ is not finitely real-EML representable.
\end{corollary}

\begin{remark}[Complex bypass]
  The barrier is specific to the \emph{real} grammar.  Over~$\C$, one has
  $\sin(x) = \mathrm{Im}(e^{ix}) = \mathrm{Im}\bigl(\eml(ix, 0)\bigr)$ for
  formal complex inputs, and the Euler identity $\sin(x) = (e^{ix} - e^{-ix})/(2i)$
  gives a two-node complex EML tree (depth 1 under a suitable complex grammar).
  The Infinite Zeros Barrier is therefore a statement about the \emph{real}
  domain only; it does not obstruct complex-EML representability.
\end{remark}

% ─────────────────────────────────────────────────────────────────────────────
\section{Computational Confirmation}
% ─────────────────────────────────────────────────────────────────────────────

To complement the analytical proof, we ran an exhaustive search over all real
EML trees with at most~$N$ internal nodes.

\begin{table}[h]
\centering
\caption{Exhaustive search results: best MSE against $\sin(x)$ over 8 probe
  points in $[0.1, 1.7]$, across all EML trees with $\leq N$ internal nodes
  and leaves in $\{x, 1\}$.}
\label{tab:search}
\begin{tabular}{rrrr}
\toprule
$N$ & Trees searched & Best MSE & Time (sec) \\
\midrule
5  & 428         & $>10^{-1}$ & $<0.01$ \\
7  & 4,862       & $>10^{-1}$ & $<0.01$ \\
9  & 58,786      & $>10^{-1}$ & 0.3 \\
11 & 742,900     & $>10^{-1}$ & 8.1 \\
12 & 1,704,034,304 & $1.12 \times 10^{-1}$ & 224.7 \\
\bottomrule
\end{tabular}
\end{table}

No tree with $N \leq 11$ internal nodes achieves MSE below $10^{-1}$ against
$\sin(x)$.  The N=12 search is ongoing; the proof guarantees the result
regardless of its outcome.

Source code: \url{https://github.com/agent-maestro/monogate},
directory \texttt{monogate-core/src/bin/n12\_search.rs}.

% ─────────────────────────────────────────────────────────────────────────────
\section*{Acknowledgements}
% ─────────────────────────────────────────────────────────────────────────────

The EML operator was introduced by Odrzywołek~\cite{odzryzwolek2026eml}.
This note proves a negative result about real representability that was left
open in that work.

% ─────────────────────────────────────────────────────────────────────────────
\begin{thebibliography}{9}

\bibitem{odzryzwolek2026eml}
A.~Odrzywołek,
\emph{The EML Operator: A Universal Generator for Elementary Functions},
arXiv:2603.21852, 2026.

\bibitem{krantz2002primer}
S.~G.~Krantz and H.~R.~Parks,
\emph{A Primer of Real Analytic Functions}, 2nd ed.,
Birkh\"{a}user, Boston, 2002.

\end{thebibliography}

\end{document}
